\section*{4. 与线束相伴的爆破的上同调}

\subsection*{4.1.}
我们回到 3.1 的情形, 它可用交换图概括如下 (见 2.1.6.3):
\begin{equation}\tag{4.1.1}
\resizebox{\textwidth}{!}{%
\begin{tikzcd}[ampersand replacement=\&,row sep=huge,column sep=huge]
X\times\mathbb P^1
  \& \widetilde X
      \arrow[l,hook',"f\times\rho"]
      \arrow[d,"f"]
      \arrow[dl,"\rho"']
  \& \widetilde\Delta
      \arrow[l,hook',"i"]
      \arrow[d,"g"]
      \arrow[r,"{(g,\rho_i)}","\sim"']
  \& \Delta\times\mathbb P^1 \\
\mathbb P^1 \& X \& \Delta\arrow[l,hook',"j"] \& {}
\end{tikzcd}}
\end{equation}
我们先在此情形下将 (2.2.2) 写得更明确.

\subsection*{命题 4.2.}
\underline{在 3.1 的假设下:}

\smallskip
\noindent
\textup{(4.2.1)} \underline{$\mathbb Z_\ell$-模的同态}
\[
\begin{tikzcd}[column sep=huge]
H^\bullet(\widetilde X,\mathbb Z_\ell)
  \arrow[r,shift left=.7ex,"f_*\oplus(1-f^*f_*)"]
& H^\bullet(X,\mathbb Z_\ell)\oplus
   \operatorname{Ker}\!\left(f_*:H^\bullet(\widetilde X,\mathbb Z_\ell)
       \to H^\bullet(X,\mathbb Z_\ell)\right)
  \arrow[l,shift left=.7ex,"f^*+\operatorname{incl}"']
\end{tikzcd}
\]
\underline{为互逆的同构.}

\smallskip
\noindent
\textup{(4.2.2)} \underline{$\mathbb Z_\ell$-模的同态}
\[
\begin{tikzcd}[column sep=huge]
H^\bullet(\widetilde X,\mathbb Z_\ell)
  \arrow[r,shift left=.7ex,"f_*\oplus(-g_*i^*)"]
& H^\bullet(X,\mathbb Z_\ell)\oplus
   H^{\bullet-2}(\Delta,\mathbb Z_\ell(-1))
  \arrow[l,shift left=.7ex,"f^*+i_*g^*"']
\end{tikzcd}
\]
\underline{为互逆的同构.}\footnote{印本在 (4.2.2)、(4.2.6)、(4.2.9) 和 (4.2.10) 中略去了 $g_*i^*$ 前的负号. 但引理 4.3 给出 $g_*i^*i_*g^*=-\mathrm{id}$, 而印本中的乘积计算本身也用了 $-g_*i^*$. 因此统一恢复此符号.}

\smallskip
\noindent
\textup{(4.2.3)} \underline{在下式上的代数结构}
\[
H^\bullet(X,\mathbb Z_\ell)\oplus
H^{\bullet-2}(\Delta,\mathbb Z_\ell(-1)),
\]
\underline{由 (4.2.2) 从下式上的杯积传递而来:}
$H^\bullet(\widetilde X,\mathbb Z_\ell)$,
\underline{对}
\[
a,b\in H^\bullet(X,\mathbb Z_\ell),
\qquad
x,y\in H^{\bullet-2}(\Delta,\mathbb Z_\ell(-1)),
\]
\underline{由下列公式表示}
\begin{equation}\tag{4.2.3.1}
(0\oplus x)\wedge(0\oplus y)
=-j_*(xy)\oplus2xy\,c_1(\underline{\mathcal O}_\Delta(1)),
\end{equation}
\begin{equation}\tag{4.2.3.2}
(a\oplus0)\wedge(b\oplus0)=ab\oplus0,
\end{equation}
\begin{equation}\tag{4.2.3.3}
(a\oplus0)\wedge(0\oplus y)=0\oplus j^*(a)y,
\end{equation}
\begin{equation}\tag{4.2.3.4}
(0\oplus x)\wedge(b\oplus0)=0\oplus xj^*(b).
\end{equation}
\underline{特别地, 若}
\[
a\oplus x\in H^\bullet(X,\mathbb Z_\ell)
  \oplus H^{\bullet-2}(\Delta,\mathbb Z_\ell(-1))
\]
以及
\[
b\oplus y\in H^{2n-\bullet}(X,\mathbb Z_\ell)
  \oplus H^{2n-2-\bullet}(\Delta,\mathbb Z_\ell(-1)),
\]
\underline{则}
\begin{equation}\tag{4.2.3.5}
(a\oplus x)\wedge(b\oplus y)=ab-j_*(xy).
\end{equation}
换言之, $H^\bullet(\widetilde X,\mathbb Z_\ell)$ 的总相交矩阵为
\[
\left(
\begin{array}{c|c}
\begin{array}{c}
\text{总相交}\\[-1pt]
\text{矩阵 ($X$ 上)}
\end{array} & 0\\ \hline
0 &
-\begin{array}{c}
\text{总相交}\\[-1pt]
\text{矩阵 ($\Delta$ 上)}
\end{array}
\end{array}
\right).
\]

\underline{证明.}
我们已经证明了 (4.2.1); 见 (2.2.22). 对于 (4.2.2), 首先注意, 由于 $g:\widetilde\Delta\to\Delta$ 为射影直线丛, 由 (1.2.2), 序列
\begin{equation}\tag{4.2.4}
0\longrightarrow H^\bullet(\Delta,\mathbb Z_\ell)
\xrightarrow{g^*}H^\bullet(\widetilde\Delta,\mathbb Z_\ell)
\xrightarrow{g_*}H^{\bullet-2}(\Delta,\mathbb Z_\ell(-1))
\longrightarrow0
\end{equation}
正合. 结合 (2.2.26), 得到如下正合列的交换图
\begin{equation}\tag{4.2.5}
\resizebox{\textwidth}{!}{%
\begin{tikzcd}[ampersand replacement=\&,column sep=large,row sep=huge]
0\arrow[r]
  \&H^\bullet(X,\mathbb Z_\ell)
       \arrow[r,"f^*"]\arrow[d,"j^*" right]
  \&H^\bullet(\widetilde X,\mathbb Z_\ell)
       \arrow[r,"g_*i^*"]\arrow[d,"i^*" right]
  \&H^{\bullet-2}(\Delta,\mathbb Z_\ell(-1))
       \arrow[r]\arrow[d,equal]
  \&0\\
0\arrow[r]
  \&H^\bullet(\Delta,\mathbb Z_\ell)
       \arrow[r,"g^*"]
  \&H^\bullet(\widetilde\Delta,\mathbb Z_\ell)
       \arrow[r,"g_*"]
  \&H^{\bullet-2}(\Delta,\mathbb Z_\ell(-1))
       \arrow[r]
  \&0.
\end{tikzcd}}
\end{equation}
由分裂 (4.2.1), 得同构
\begin{equation}\tag{4.2.6}
f_*\oplus(-g_*i^*):
H^\bullet(\widetilde X,\mathbb Z_\ell)
\xrightarrow{\sim}
H^\bullet(X,\mathbb Z_\ell)\oplus
H^{\bullet-2}(\Delta,\mathbb Z_\ell(-1)).
\end{equation}
将 (4.2.5) 转置, 并反转次数 $\bullet\mapsto2n-\bullet$, 得到正合列的交换图
\begin{equation}\tag{4.2.7}
\resizebox{\textwidth}{!}{%
\begin{tikzcd}[ampersand replacement=\&,column sep=large,row sep=large]
0
  \&H^\bullet(X,\mathbb Z_\ell)\arrow[l]
  \&H^\bullet(\widetilde X,\mathbb Z_\ell)\arrow[l,"f_*"']
  \&H^{\bullet-2}(\Delta,\mathbb Z_\ell(-1))
       \arrow[l,"i_*g^*"']
  \&0\arrow[l]\\
0
  \&H^{\bullet-4}(\Delta,\mathbb Z_\ell(-2))
       \arrow[l]\arrow[u,"j_*"']
  \&H^{\bullet-2}(\widetilde\Delta,\mathbb Z_\ell(-1))
       \arrow[l,"g_*"']\arrow[u,"i_*"']
  \&H^{\bullet-2}(\Delta,\mathbb Z_\ell(-1))
       \arrow[l,"g^*"']\arrow[u,equal]
  \&0\arrow[l]
\end{tikzcd}}
\end{equation}
\begin{editorialnote}
原印本及英译底本在图 (4.2.7) 左下写作 $H^{\bullet-2}(\Delta,\mathbb Z_\ell(-2))$, 并将右下方的水平箭头标为 $g_*$. 由于 $j$ 的余维为二, $j_*$ 将次数增加四; 由于 $g$ 的相对维数为一, $g_*$ 将次数降低二. 因而左下次数应为 $\bullet-4$. 右下箭头从 $\Delta$ 指向 $\widetilde\Delta$, 应为 $g^*$. 这也正是 (4.2.5) 在 Poincar\'e 对偶下的转置. 此处据此订正这两处, 未改变其余节点或箭头.
\end{editorialnote}

从而得到同构
\begin{equation}\tag{4.2.8}
H^\bullet(X,\mathbb Z_\ell)\oplus
H^{\bullet-2}(\Delta,\mathbb Z_\ell(-1))
\xrightarrow[\sim]{\ f^*+i_*g^*\ }
H^\bullet(\widetilde X,\mathbb Z_\ell).
\end{equation}
